\chapterimage{figs/chapterhead.png}
\chapter{Model Language for Users}
\label{chap:model-language-for-users}

This chapter documents the current user-facing textual representation of a
GenESyS model. It does not try to replace the graphical workflow or the later
developer-level language reference. Instead, it explains how the textual model
view fits into the GUI, how to read a saved model file, and how to keep the
text synchronized with the current kernel model.

\noindent\textbf{Objective.} Show how the textual model language relates to the
graphical model and how a user can read it as a second view of the same system.

\noindent\textbf{Audience.} Users who already know the GUI and want to inspect
or edit the textual model representation with confidence.

\noindent\textbf{Scope.} Cover the current text editor surface, the current
line-oriented file structure, the relationship between model text and GUI
elements, and the practical limits of this chapter.

\noindent\textbf{Status.} Partially validated by code inspection and example
model files. The chapter matches the current GUI services
(\texttt{ModelLanguageSynchronizer} and \texttt{GraphicalModelSerializer}) and
the current \texttt{.gen} file shape, but it intentionally stops short of a
full grammar reference.

The chapter follows this sequence:
Figure~\ref{fig:model-language-tab},
Table~\ref{tab:model-language-sections}, and
Figure~\ref{fig:model-language-flow}.

\section{What the model language is}
\label{sec:model-language-what-it-is}

The user-facing model language is the textual form of a GenESyS model. It is
the same model that appears in the GUI, but serialized as text so it can be
read, saved, copied, compared, and in some cases edited directly.

In the current code, the GUI keeps this representation in the
\texttt{TextCodeEditor} widget inside the model-language tab set. The
\texttt{ModelLanguageSynchronizer} service rebuilds that text from the current
\texttt{Simulator} state, and it can also apply textual edits back into the
model. The persistence path is handled by
\texttt{GraphicalModelSerializer}, which saves the textual representation using
the established compatibility format.

That means the text is not a separate abstraction with its own independent
truth. It is a synchronized view of the current model. If the model changes in
the GUI, the text view must be refreshed. If the text is edited, the model must
be revalidated before it can be used again.

\begin{figure}[htbp]
\centering
% FIGURE-SPEC-BEGIN
% ID: fig:model-language-tab
% Source of truth:
%   - source/applications/gui/genesys/mainwindow.ui
%   - source/applications/gui/genesys/mainwindow.cpp
%   - source/applications/gui/genesys/services/ModelLanguageSynchronizer.h
%   - source/applications/gui/genesys/services/GraphicalModelSerializer.h
%   - models/AirportSecurityExample.gen
% Purpose:
%   Show the current model-language editor/tab and a representative text excerpt.
% Required visual content:
%   - the model-language tab or editor surface
%   - a real model text excerpt
%   - a visible connection to the surrounding GUI
% Accessibility description:
%   A screenshot that makes it clear that the text is the model's serialized view, not a separate document.
% Validation:
%   The visible widget labels and the excerpt structure must match the current repository code and saved model files.
% FIGURE-SPEC-END
\figureplaceholder{figs/generated/model-language-tab}
\caption{Textual model view used by the GUI to represent the current model.}
\label{fig:model-language-tab}
\end{figure}

Figure~\ref{fig:model-language-tab} should be read as a direct companion to the
graphical model, not as a replacement for it.

\section{How a saved model is organized}
\label{sec:model-language-organization}

A current saved model file is line-oriented and sectioned. One current example
is the airport security example model under \texttt{models/}. Its
\texttt{.gen} file starts with comment headers, then lists the simulator and
model metadata, and then continues with data definitions and model components.

The exact naming of sections can vary slightly between models, but the current
shape is easy to recognize:
\begin{itemize}
\item comment lines at the top identify the file as a GenESyS model;
\item a header section records simulator, model, and simulation metadata;
\item model data definitions appear next;
\item model components follow after that.
\end{itemize}

The format is intentionally explicit. A line typically begins with a numeric
prefix, then a type name, then a quoted object name, and then a list of
property assignments. For example, the example model file contains entries such
as \texttt{Simulator}, \texttt{ModelInfo}, \texttt{ModelSimulation},
\texttt{EntityType}, \texttt{Resource}, \texttt{Queue}, \texttt{Create},
\texttt{Process}, \texttt{Seize}, \texttt{Delay}, \texttt{Release},
\texttt{Decide}, and \texttt{Dispose}.

\begin{table}[htbp]
\centering
\caption{Typical sections of a current GenESyS model text file.}
\label{tab:model-language-sections}
\small
\begin{tabular}{|p{3.5cm}|p{7.4cm}|}
\hline
\textbf{Section} & \textbf{What the user reads there} \\
\hline
Comment header & File identity and simple human-readable reminders \\
\hline
Simulator and model metadata & Names, descriptions, version metadata, and simulation defaults \\
\hline
Model data definitions & Entity types, attributes, queues, resources, and other supporting data \\
\hline
Model components & The connected flow elements that form the actual model behavior \\
\hline
\end{tabular}
\end{table}

The model text in Table~\ref{tab:model-language-sections} is useful because it
maps directly to the same conceptual objects the user already sees in the GUI.
That one-to-one relation is the main reason the textual representation is worth
learning.

\section{How to read the text as a user}
\label{sec:model-language-how-to-read}

The best reading strategy is to start from the same questions used in the GUI.
Instead of trying to understand the whole file at once, ask what each section
describes in the graphical model.

In practice:
\begin{itemize}
\item the header tells you which model is open and which simulation defaults it uses;
\item data definitions tell you what reusable objects the model depends on;
\item component lines tell you how the flow is wired;
\item expressions tell you what behavior is parameterized rather than hard-coded.
\end{itemize}

That same example model is a useful reading aid because it is small enough to
inspect and still contains all the major current categories. Its first
component block shows how text expresses the same model that the GUI diagram
represents: the source component, the next component, the delay expression,
the queue relation, the resource request, the decision condition, and the
alternative exits are all present in text.

For a user, the important habit is to match names and links:
\begin{itemize}
\item component names in the file should match the labels seen in the GUI;
\item property values should match the properties edited in the property panel;
\item connection fields such as \texttt{nextId} should line up with the visible flow;
\item captions and descriptions should remain readable as plain text.
\end{itemize}

\begin{figure}[htbp]
\centering
% FIGURE-SPEC-BEGIN
% ID: fig:model-language-flow
% Source of truth:
%   - models/AirportSecurityExample.gen
%   - source/applications/gui/genesys/services/ModelLanguageSynchronizer.h
%   - source/applications/gui/genesys/services/GraphicalModelSerializer.h
% Purpose:
%   Show the practical reading flow from GUI model to text file and back again.
% Required visual content:
%   - a graphical model view
%   - a corresponding model-language excerpt
%   - arrows or annotations showing round-trip synchronization
% Accessibility description:
%   A simple explanatory diagram that helps the reader map visual objects to text lines and back.
% Validation:
%   The diagram must describe the current synchronization path, not a hypothetical one.
% FIGURE-SPEC-END
\figureplaceholder{figs/generated/model-language-flow}
\caption{How the textual model view maps to the graphical model and back.}
\label{fig:model-language-flow}
\end{figure}

Figure~\ref{fig:model-language-flow} summarizes the practical reading rule:
learn the graphical model first, then use the textual view to confirm and
compare what the GUI already showed.

\section{Editing and synchronization}
\label{sec:model-language-editing}

The current implementation supports a synchronization workflow, not a free-form
text editor detached from the simulator. That distinction matters.

When the text is changed, the GUI must apply the text back into the current
model and then refresh the rest of the interface. When the model changes in the
GUI, the text must be regenerated so the editor does not go stale. The code
paths involved in this cycle are already explicit in the repository:
\begin{itemize}
\item The sync service rebuilds the text from model state.
\item The sync service applies text edits back to the model.
\item The serializer writes the current text representation to disk.
\item The serializer restores persisted model state.
\end{itemize}

For the user, the conservative rule is simple: after a direct text change,
reload or re-sync the model before relying on the GUI again. After a GUI change,
confirm that the text view has been refreshed before assuming the file on disk
already reflects the new state.

This chapter does not introduce a separate command line or a new syntax for
manual editing. It documents the existing synchronized surface only.

\section{What this chapter does not cover}
\label{sec:model-language-limits}

This chapter is intentionally not the full language reference. It does not try
to define the complete grammar, every expression function, or every parser
production. Those topics belong to the later developer-level language chapter.

It also does not claim that every old or future file variant is supported
forever. The only safe claim here is about the current repository state: the
GUI stores a textual model view, saves it in the established format, and keeps
that view synchronized with the kernel model.

If a text excerpt and the GUI appear to disagree, the user should treat that as
a synchronization problem to resolve, not as a signal that the text is a better
authority than the model or vice versa.

\section{Final considerations}
\label{sec:model-language-final-considerations}

The main value of the model language for users is readability. It gives the
same model a second, formal surface that can be inspected alongside the GUI.
That helps with understanding, comparison, review, and troubleshooting, but it
only works well when the user remembers that the text and the diagram are two
views of one synchronized model.

In the reading sequence used by this manual, the textual model view comes after
the GUI, the first model, editing, execution, and results. That order is
intentional. By the time the reader reaches this chapter, the model language is
no longer an opaque notation. It is a familiar model written in another form.

The next chapter returns to command-line usage for users who need the terminal
workflow or want a more scripted interaction style.
